\chapter{Historial de cambios de superficie}
\label{app:spec-changelog}

\section{Política de versiones}

\textbf{HTTP API} y \textbf{MCP tools} tienen ciclo independiente.
Cada una sigue semver:

\begin{itemize}
    \item \textbf{Mayor}: cambio incompatible (remoción de
    endpoint/tool, cambio de shape).
    \item \textbf{Menor}: agregado compatible (nuevo endpoint, campo
    opcional).
    \item \textbf{Patch}: corrección sin cambio de shape.
\end{itemize}

El DSL tiene su propio versionado (\code{v1} estable, \code{v2}
experimental detrás de feature flag). La C ABI tiene versión
entera monotónica (\code{gm_abi_version()}).

\section{v1.0 — baseline estable}

Corresponde al cierre del rediseño shearing-layers. Lo que entra a
v1.0:

\textbf{HTTP API}: 62 endpoints de producto + \code{/} + \code{/metrics}.
Todos los listados en cap.~\ref{sec:api-http} son \code{stable}
salvo:

\begin{itemize}
    \item \code{POST /publish_iocs_to_sentinel} — \code{experimental}
    (gated por feature flag).
    \item \code{POST /file_ticket} — \code{experimental} (gated por
    feature flag).
    \item \code{POST /kql} — \code{experimental} (depende de credenciales
    externas).
\end{itemize}

\textbf{MCP tools}: 55 tools en 10 dominios. Todas \code{stable}
salvo:

\begin{itemize}
    \item \code{integrations.publish_iocs_to_sentinel}.
    \item \code{integrations.file_ticket}.
    \item \code{scoring.enable_anomaly_scoring} cuando
    \code{w5_gnn_threat > 0}.
\end{itemize}

\textbf{DSL}: gramática v1 (ver cap.~\ref{sec:dsl}).

\textbf{C ABI}: \code{gm_abi_version() = 1}.

\section{Compromisos v1}

\begin{itemize}
    \item Un endpoint \code{stable} no se elimina ni cambia de
    método/path dentro de v1. Se puede agregar campo opcional nuevo
    al request o response.
    \item Una tool MCP \code{stable} no cambia su \code{inputSchema}
    de forma incompatible dentro de v1.
    \item Una hipótesis DSL válida en v1 sigue parseando igual en
    toda la rama v1.
    \item \code{gm_abi_version()} cambia si y sólo si hay cambio
    incompatible en la C ABI.
\end{itemize}

\section{Qué queda fuera de v1}

\begin{itemize}
    \item \textbf{Agentic slow-lane} — \code{agentic_*} endpoints y
    tools. Interfaces funcionan pero se considera
    \code{experimental} mientras no se cierre el backend de LLM
    local.
    \item \textbf{GNN v2} — model gate y
    \code{scoring.explain_score.threat_class} sólo son relevantes si
    el usuario provee modelo ONNX. Sin modelo, comportamiento es
    estable (devuelve 0.0).
    \item \textbf{DSL v2} — operadores temporales explícitos
    (\code{before}, \code{within}) detrás del feature
    \code{dsl-v2-experimental}.
\end{itemize}

\section{Tags git de release}

\begin{itemize}
    \item \code{platform/api-v1.0.0} — primer tag estable de la API
    HTTP.
    \item \code{platform/mcp-v1.0.0} — primer tag estable del
    servidor MCP.
    \item \code{libgraphmatch-v1.0.0} — C ABI congelada en
    \code{gm_abi_version() = 1}.
\end{itemize}

El workspace raíz sigue semver de producto (versión del desktop
app), independiente de estos.

\section{Fuente canónica}

El listado \emph{vivo} de endpoints es \code{GET /v1/schema}. Si hay
discrepancia entre este documento y la respuesta del daemon, el
daemon gana; este documento se actualiza en el siguiente release.

Para tools MCP, el listado canónico es
\code{GET /v1/mcp/schema} (emitido por el servidor MCP) o
\filepath{platform/mcp/src/registry.ts} en el código fuente.
